% ============================================================================
% MODULE 7: CLIMATE SCENARIOS & FUTURE OUTLOOK
% ============================================================================
% Duration: 60 minutes (adaptation, climate change, decision frameworks)
% Learning Goals: Apply APSIM to future climate; identify adaptation strategies
% Prerequisites: All previous modules

\chapter{Module 7: Climate Scenarios \& Future Outlook}
\label{ch:module7}

\section{Learning Objectives}

By the end of this module, you will:
\begin{itemize}
  \item Understand how to construct future climate scenarios
  \item Apply APSIM under baseline vs. climate change weather
  \item Identify crop varieties and management options that maintain yield
  \item Assess adaptation feasibility and economics
  \item Communicate uncertainty to farmers/policymakers
\end{itemize}

\section{Climate Change and Agriculture}

\subsection{Global Context}

\begin{itemize}
  \item Temperature rising: ~1.1°C above pre-industrial (as of 2024)
  \item Rainfall patterns shifting: Some regions wetter, others drier
  \item Extreme events increasing: Droughts, floods, heat waves
  \item Crop modeling crucial for adaptation planning
\end{itemize}

\section{Constructing Climate Scenarios}

\subsection{GCM (Global Climate Model) Outputs}

Climate models (e.g., IPCC scenarios) provide future temperature and rainfall projections:
\begin{itemize}
  \item \textbf{RCP / SSP scenarios:} Different emission pathways (low, medium, high)
  \item \textbf{Spatial resolution:} Typically 1--50 km grid cells
  \item \textbf{Uncertainty:} Multiple GCMs, ensemble averages
\end{itemize}

\subsection{Downscaling to Farm Scale}

Global Climate Models (GCMs) operate at coarse spatial resolution (50--200 km grids). Before using GCM outputs in APSIM, projections must be \emph{downscaled} to the local scale. Common approaches:

\begin{enumerate}
  \item \textbf{Delta method (simplest)}: Add a monthly temperature offset and multiply rainfall by a scaling factor derived from GCM projections.
    \begin{itemize}
      \item Example: GCM predicts +1.5\textdegree C warming by 2050 $\to$ add 1.5 to all T$_{max}$ and T$_{min}$ values in the historical weather file
      \item For rainfall: GCM predicts $-$10\% winter rainfall $\to$ multiply monthly rainfall by 0.90
    \end{itemize}
  \item \textbf{Bias correction}: Statistical adjustment so the modelled baseline matches observed historical data (removes systematic GCM biases)
  \item \textbf{Statistical downscaling}: Relate large-scale atmospheric patterns to local precipitation using regression or weather generators
\end{enumerate}

\begin{tipbox}
\textbf{For this tutorial}, if you construct a climate scenario, use the delta method: it is transparent, easy to interpret, and sufficient for exploring sensitivity. Open the \filename{.met} weather file in a text editor and modify the temperature columns.
\end{tipbox}

\section{Adaptation Strategies}

\subsection{Genetic (Variety Selection)}

\begin{itemize}
  \item Shift to heat-tolerant varieties
  \item Shift to drought-tolerant varieties (deep roots, higher transpiration efficiency)
  \item Shift to earlier-maturing varieties (escape late-season drought)
\end{itemize}

\subsection{Management (Practice Changes)}

\begin{itemize}
  \item Earlier sowing (cooler temperatures pre-anthesis)
  \item Increased N to maintain protein in warmer years
  \item Soil water conservation (mulch, reduced tillage)
  \item Diversification (intercropping, rotation)
\end{itemize}

\subsection{Infrastructure (Structural Adaptation)}

\begin{itemize}
  \item Irrigation (if water available)
  \item Soil improvement (increase water-holding capacity)
  \item Crop insurance (financial risk transfer)
\end{itemize}

\section{APSIM-Based Adaptation Assessment}

\subsection{Step-by-Step: Constructing a Climate Scenario in APSIM}

\begin{enumerate}
  \item \textbf{Establish a baseline}: Run your 20-year simulation with historical weather (already done in Exercise D). Record mean yield, CV, and P(failure).

  \item \textbf{Create a modified weather file}: Make a copy of the \filename{.met} file. Apply the delta method:
    \begin{itemize}
      \item Add +1.5\textdegree C to all T$_{max}$ and T$_{min}$ values (approximate 2050 projection for subtropical QLD)
      \item Reduce in-season (May--September) rainfall by 10\%
    \end{itemize}

  \item \textbf{Update the Weather component}: In APSIM, click the \apsimterm{Weather} component $\to$ change \variable{FileName} to your modified weather file.

  \item \textbf{Re-run the simulation}: Keep all other settings identical.

  \item \textbf{Compare results}: In Excel, place baseline and scenario yields side by side. Calculate:
    \begin{itemize}
      \item Change in mean yield (kg/ha and \%)
      \item Change in yield CV
      \item Change in P(failure)
    \end{itemize}
\end{enumerate}

% ---- IMAGE PLACEHOLDER ----
\begin{figure}[H]
  \centering
  \fbox{\parbox{0.85\textwidth}{\centering\vspace{1.5cm}
    \textit{[IMAGE: Side-by-side bar chart comparing annual yields under Baseline vs. +1.5\textdegree C scenario.}\\
    \textit{Annotate: mean yield for each scenario as a horizontal dashed line; label years where scenario is notably lower.]}
  \vspace{1.5cm}}}
  \caption{Comparison of simulated wheat yields under historical baseline and a +1.5\textdegree C warming scenario. Increased temperature accelerates development and shortens grain fill duration.}
  \label{fig:climate-scenario}
\end{figure}
% ---- END PLACEHOLDER ----

\subsection{Illustrative Results}

Indicative outcomes for a +1.5\textdegree C / $-$10\% rainfall scenario at Northam (WA Wheatbelt):

\begin{center}
\begin{tabular}{l|r|r|r}
\textbf{Metric} & \textbf{Baseline} & \textbf{+1.5\textdegree C / $-$10\% rain} & \textbf{Change} \\
\hline
Mean yield (kg/ha) & 2800 & 2350 & $-$16\% \\
Yield CV (\%) & 28 & 36 & +8 points \\
P(yield $<$ 1500 kg/ha) & 15\% & 28\% & +13 points \\
Season length (days) & 195 & 180 & $-$15 days \\
\end{tabular}
\end{center}

(Indicative values for illustration only.)

\begin{keyconceptbox}
\textbf{Two effects of warming in APSIM:}
\begin{enumerate}
  \item \textbf{Shorter season}: More GDD/day means the crop reaches maturity faster, reducing grain fill duration
  \item \textbf{Increased water demand}: Higher temperature raises PET, depleting ESW faster during grain fill
\end{enumerate}
Both effects reduce yield in most scenarios. Adaptation strategies target both.
\end{keyconceptbox}

\section{Communicating Uncertainty}

\begin{warningbox}
\textbf{Challenge:} APSIM forecasts have multiple sources of uncertainty:
\begin{itemize}
  \item Climate model uncertainty (different GCMs give different projections)
  \item Downscaling uncertainty (statistical method choices)
  \item Crop model parameter uncertainty (calibration, regional differences)
  \item Natural variability (internal climate noise; chaotic system)
\end{itemize}

\textit{Solution:} Present ranges (percentiles), not point estimates. Use ensemble approaches (multiple models, multiple years).
\end{warningbox}

\section{Decision Frameworks}

\subsection{Robust Decision Making}

Choose strategies that perform reasonably well across multiple scenarios (not optimized for one scenario only).

Example: N rate of 120 kg/ha may not be optimal in any single year, but performs well across dry, normal, and wet years.

\section{Sustainability Considerations}

Climate adaptation must also be sustainable:

\begin{itemize}
  \item \textbf{N leaching}: Higher rainfall intensity under some climate scenarios increases drainage and N leaching. APSIM's NO$_3^-$ output tracks this; excess leaching is both an economic loss and an environmental risk.
  \item \textbf{Water use}: Irrigation as an adaptation strategy requires water availability, which may also be declining in some regions.
  \item \textbf{Soil carbon}: Management changes (tillage, residue) affect soil organic matter, which feeds back to mineralisation rates and water-holding capacity --- both captured in APSIM's SoilOrganicMatter component.
  \item \textbf{Social acceptance}: Variety changes and new practices require knowledge transfer, seed supply chains, and economic support for adoption.
\end{itemize}

\begin{tipbox}
\textbf{Triple bottom line:} The best adaptation strategies improve or maintain yield (economic), reduce environmental impact (ecological), and are feasible for farmers to adopt (social). Use APSIM to assess the first two; consult socioeconomic literature for the third.
\end{tipbox}

\section{Summary}

\begin{keyconceptbox}
\textbf{Module 7 key points:}
\begin{enumerate}
  \item GCM outputs must be downscaled to farm scale before use in APSIM; the delta method (add T offset, multiply rainfall) is the simplest approach
  \item Warming shortens season length (faster GDD accumulation) and increases water demand (higher PET)
  \item Multiple RCP/SSP scenarios produce a range of projections --- present as ranges, not single values
  \item Adaptation strategies: genetic (heat-tolerant varieties), management (earlier sowing, split N), infrastructure (irrigation)
  \item Robust decision making: choose strategies that perform acceptably across multiple scenarios
  \item Climate adaptation must also be sustainable --- consider N leaching, water use, and social feasibility
\end{enumerate}
\end{keyconceptbox}

\section{Synthesis: APSIM in Your Career}

This tutorial has covered:
\begin{itemize}
  \item Simulation mechanics (Module 1)
  \item Hands-on practice (Modules 2--3, Exercises A--D)
  \item Agronomic applications (Modules 4--5)
  \item Risk and adaptation (Modules 6--7)
\end{itemize}

APSIM is now a tool in your toolkit for:
\begin{itemize}
  \item Research questions (hypothesis testing via simulation)
  \item Extension advice (quantifying management options)
  \item Operational decisions (sowing timing, N optimization)
  \item Policy analysis (climate impact, food security)
\end{itemize}

\section{Continuing Your Learning}

\begin{itemize}
  \item APSIM documentation: \url{https://www.apsim.info}
  \item Community forum: APSIM Next Gen GitHub discussions
  \item Advanced training: APSIM workshops (CIMMYT, ICRISAT, regional centers)
  \item Job opportunities: APSIM skills highly valued in ag tech, research, extension
\end{itemize}

% ============================================================================
% END MODULE 7
% ============================================================================
